% TOP_PROBLEM: {"id": 2785, "title": "Kirby Problem 2.37", "category_id": 7, "category": {"id": 7, "name": "topology", "display_name": "Topology", "description": "Properties preserved under continuous deformations.", "slug": "topology", "order_index": 7, "created_at": "2026-07-31T15:26:25.671Z"}, "status": "open"}
% FIRST_POSTED: null
% ENTRY_KIND: statement_only
% Imported from: batch11.tex; UnsolvedMath ID 2785
\documentclass[11pt]{article}
\usepackage[margin=25mm]{geometry}
\usepackage{fontspec}
\setmainfont{Latin Modern Roman}
\setsansfont{Latin Modern Sans}
\usepackage{amsmath,amssymb,mathtools}
\usepackage{unicode-math}
\setmathfont{Latin Modern Math}
\usepackage{xurl,hyperref}
\setcounter{secnumdepth}{0}
\setlength{\parindent}{0pt}
\setlength{\parskip}{5pt}
\setlength{\emergencystretch}{3em}
\newcommand{\sref}[2]{\hyperlink{source:#1:#2}{[S#2]}}
\newcommand{\cataloguescope}[1]{\paragraph{Recorded target.}#1}
\begin{document}
% UnsolvedMath ID 2785; source catalogue code KP-2.37
\setcounter{section}{2784}
\section{Kirby Problem 2.37}\label{2785}
{\small\sffamily UnsolvedMath ID 2785 · Topology}
\subsection{Definitions and mathematical statement}

\paragraph{Shared notation.}$\mathbb N=\{1,2,\ldots\}$, and $\mathbb Z,\mathbb Q,\mathbb R,\mathbb C$ denote the integers, rationals, reals and complex numbers. Any other conventions are specified locally.


An oriented surface $S$ is of infinite type if it is not a compact finite-genus surface with finitely many punctures and boundary circles. Its mapping class group $\operatorname{Mod}(S)$ consists of orientation-preserving homeomorphisms modulo isotopy, fixing boundary pointwise. An end is an element of the inverse limit of the sets of components of complements of compact subsets. An end is planar if it has a genus-zero neighborhood; otherwise it is accumulated by genus. For finite-type surfaces, a mapping class is periodic when it has finite order, reducible when it preserves an essential finite multicurve, and otherwise pseudo-Anosov: a representative preserves transverse measured foliations while expanding one and contracting the other by a factor $\lambda>1$. An essential curve is neither null-homotopic nor peripheral.
\paragraph{Statement recorded in the supplied source.}
Develop an exhaustive Nielsen--Thurston-type classification for mapping classes of arbitrary infinite-type surfaces. Determine the appropriate replacement for reducibility and identify the classes of homeomorphisms that generalize pseudo-Anosov behavior, allowing the additional categories required in infinite type. The classification is requested for all mapping classes, not only tame or end-periodic ones. \sref{2785}{2}

\subsection{Short English statement}
Classify all mapping classes of infinite-type surfaces and identify which dynamical classes replace finite-type pseudo-Anosov maps.
\subsection{Sources}
\begin{description}
\item[\hypertarget{source:2785:1}{S1}] ULAM.AI and the UnsolvedMath Contributors, UnsolvedMath: A Curated Collection of Open Mathematics Problems, record 2785 (KP-2.37). CC BY 4.0. \url{https://huggingface.co/datasets/ulamai/UnsolvedMath}.
\item[\hypertarget{source:2785:2}{S2}] R. İnanç Baykur, Robion C. Kirby and Daniel Ruberman, eds., \emph{K3: A New Problem List in Low-Dimensional Topology}, Mathematical Surveys and Monographs 295, American Mathematical Society (2026), Problem 2.37 and its remarks; Chapters 1 and 2 for the stated conventions. \url{https://math.berkeley.edu/sites/default/files/surv-295-ruberman-watermarked-author-pdf.pdf}.
\end{description}
\label{end:2785}
\end{document}
